\section{Background}
% ============================================================

Matrix differential equations arise naturally in numerical linear
algebra, continuous-time optimisation, and dynamical formulations of
matrix factorisation problems. In this project, we study a nonlinear
matrix flow that continuously transforms a nonsingular matrix toward
the orthogonal factor in its polar decomposition.

For a nonsingular matrix $\mathbf{X}_0$, the polar decomposition writes
\[
\mathbf{X}_0 = \mathbf{Q}\mathbf{H},
\]
where $\mathbf{Q}$ is orthogonal and $\mathbf{H}$ is symmetric positive
definite. The orthogonal factor is important in many problems involving
matrix normalisation, optimisation on matrix spaces, and numerical
linear algebra.

The model considered in this report is the matrix differential equation
\begin{equation}
\label{eq:background_matrix_flow}
\dot{\mathbf{X}}(t)
=
\mathbf{X}(t)
\left(
\mathbf{I}_n-\mathbf{X}(t)^T\mathbf{X}(t)
\right),
\qquad
\mathbf{X}(0)=\mathbf{X}_0.
\end{equation}
For a nonsingular initial matrix, the flow preserves the singular
vectors while driving the singular values toward one. As a result,
$\mathbf{X}(t)$ approaches the orthogonal polar factor as
$t\rightarrow\infty$.

The system can also be interpreted as a gradient flow. Defining
\begin{equation}
\label{eq:background_lyapunov}
\Phi(\mathbf{X})
=
\frac14
\left\|
\mathbf{X}^T\mathbf{X}-\mathbf{I}_n
\right\|_F^2,
\end{equation}
the evolution satisfies
\[
\frac{d}{dt}\Phi(\mathbf{X}(t)) \leq 0.
\]
Therefore, the Lyapunov functional decreases along the exact trajectory.
This provides a useful structural diagnostic in addition to ordinary
solution error.

This problem is particularly suitable for numerical analysis because it
combines a nonlinear matrix-valued ODE with an exact finite-time
solution obtained from the singular value decomposition of the initial
matrix. The exact solution provides a reliable reference against which
numerical approximations can be checked.

At the same time, the model exhibits nontrivial numerical behaviour.
Large initial singular values can generate a rapid transient and impose
restrictive step sizes on explicit schemes. The problem therefore
provides a useful setting for studying the interaction between accuracy,
absolute stability, adaptive step-size control, and computational work.

In this report, Explicit Euler, classical fourth-order Runge--Kutta
(RK4), and Implicit Euler are applied to the matrix flow. Their
convergence and stability properties are analysed and tested
numerically. The computed solutions are validated using the exact SVD
solution and an independent tight-tolerance Radau reference, while
additional diagnostics include singular-value evolution, Lyapunov
decay, orthogonality defect, and rank-deficient behaviour.

% ============================================================